\documentclass[10pt,a4paper]{amsart}
\usepackage[margin=19mm]{geometry}
\usepackage{amsmath,amssymb,mathtools}
\usepackage[scr=rsfs,cal=euler]{mathalfa}
\usepackage{xfrac}
\usepackage[unicode,hidelinks,bookmarks=false]{hyperref}
\newcommand{\ie}{i.e.\ }
\newcommand{\cf}{cf.\ }
\newcommand{\resp}{respectively\ }
\newcommand{\Subsection}{\subsection}
\providecommand{\nobreakdash}{\nobreak\hbox{-}}
\setlength{\emergencystretch}{2em}
\multlinegap0pt
\usepackage{bm}
\usepackage{bbm}
\usepackage{dsfont}
\usepackage{xspace}
\usepackage{cancel}
\usepackage{mathtools}
\usepackage{amsthm}
\makeatletter
\@ifundefined{theorem}{\newtheorem{theorem}{Theorem}}{}
\makeatother
\newcommand{\F}{\mathcal{F}}
\newcommand{\Presh}{\textrm{Presh}}
\newcommand{\bool}[1]{\mathsf{#1}}
\newcommand{\bp}[1]{\left\lbrace #1 \right\rbrace}
\title[Boolean valued models, presheaves, and \'etal\'e spaces]{Boolean valued models, presheaves, and \'etal\'e spaces}
\author{Moreno Pierobon, Matteo Viale}
\date{Source: arXiv:2006.14852v6 (CC BY 4.0)}
\begin{document}
\maketitle

\noindent\emph{Source and licence.} Boolean valued models, presheaves, and \'etal\'e spaces. Version arXiv:2006.14852v6 (CC BY 4.0), distributed under the Creative Commons Attribution 4.0 International licence, as recorded in the arXiv licence metadata for this exact version and shown on the abstract page https://arxiv.org/abs/2006.14852v6. Full rights notice: licenses/arxiv-2006.14852-CC-BY-4.0.txt.\par\smallskip

\noindent\textit{Editorial scope.} Counted unit: the Theorem whose source label is \texttt{fact:lambdahaus} in the article (source0.tex lines 2685--2690, source label \texttt{fact:lambdahaus}), delivered together with the article's own complete original proof of it (source0.tex lines 2692--2762, one \texttt{proof} environment of 71 lines). No mathematics is written, repaired or re-derived: statement and proof are the article's own text line for line. Required context retained uncounted: none. Every definition, display and label of those lines is retained unchanged. Omitted from this package and not redistributed: 1--2684, 2691--2691, 2763--4918. Nothing inside a retained excerpt is omitted or altered: every retained body is delivered exactly as the article's lines are written, so no omission receipt of any kind accompanies this package. Citations: the retained excerpts carry no citation command at all, so no bibliography entry of the article is transcribed and no reference is redistributed. Reference adaptation: none was needed and none was made; every reference inside the delivered unit resolves inside it, so no unresolved reference or citation is produced. Printed slips kept literal: no printed word, symbol, hypothesis or number of the counted statement or of its proof was altered. Layout and class: the article's own class is \texttt{amsart} with packages including geometry, times, babel, inputenc, color, graphicx, fancyhdr, mathtools, csquotes, amsmath, amsthm, amssymb, amsfonts, mathrsfs; the wrapper is the lane's pinned \texttt{amsart} editorial wrapper at 10pt on A4, which restates the theorem-like environments the retained text needs and adds the article's own macro definitions that the retained text needs (\texttt{\textbackslash F}, \texttt{\textbackslash Presh}, \texttt{\textbackslash bool}, \texttt{\textbackslash bp}), with extra wrapper packages (bm, bbm, dsfont, xspace, cancel, mathtools). The delivered numbering is the unit's own. Assets: no retained span contains a figure environment, a graphics-inclusion command, a PDF-page inclusion, a table or a TikZ picture; no image file is copied into this package, the package declares no asset dependency and no image file is redistributed with it. Licence: version arXiv:2006.14852v6 of this article is distributed under the Creative Commons Attribution 4.0 International licence, as recorded in the arXiv licence metadata for this exact version and shown on the abstract page https://arxiv.org/abs/2006.14852v6. A full rights notice is delivered at licenses/arxiv-2006.14852-CC-BY-4.0.txt. Characters: every retained excerpt is pure ASCII, so the delivered engine, XeLaTeX with its default Latin Modern text font, reports no missing character.
\begin{theorem}
\label{fact:lambdahaus} \label{prop:lambdahaus}
Let $\bool{B}$ be a complete Boolean algebra and $\F\in\Presh(\bool{B})$.
The topology $\sigma_\F$ on $\Lambda^1_\F$ is Hausdorff,  locally compact, and extremally disconnected.
 %Hence $(\pi_\F, \Lambda^1_\F, \sigma_\F)$ is an \'etal\'e space.
 \end{theorem}

 \begin{proof}
\emph{}

\begin{description}
\item[$\Lambda^1_\F$ has the Hausdorff property]
The only non trivial case occurs when the two points we need to separate $[f]_{G}\neq[g]_{H}$ are on the same stalk; we need to show that these two points can be separated by basic open sets. This case occurs when $G=H$ and $f\in\F(b_f)$, $g\in\F(b_g)$ with $N_{b_f},N_{b_g}\in G$. Let%\footnote{It is here that we crucially use that the base space of the bundle $\Lambda^J_\F$ is extremally disconnected and $P=\bool{B}^+$ for some complete Boolean algebra $\bool{B}$: if $\bool{B}$ is not complete $U=\bigvee_{\RO(P)}\{N_q: f\restriction q=g\restriction q\}$ may not be in $\bool{B}$ in which case $J^1_P$ would not be defined on $U$, and one would not be able to separate $[f]_{J^1_P, G}\neq[g]_{J^1_P, G}$ using the argument below.}
 \[
 N_b:=\bigvee_{\bool{B}}\{N_q: f\restriction q=g\restriction q\}.
 \]
 We observe that $N_b\notin G$: otherwise \[N_c=N_b\cap N_{b_f}\cap N_{b_g}=\bigvee_{\bool{B}}\{N_q: f\restriction q=g\restriction q, q\leq p_f, p_g\}
 \]
 would be in $G$. This would give that $\bp{q\leq N_c: f\restriction q=g\restriction q}$ is dense below $N_c$, yielding $[f]_{G}=[g]_{G}$, a contradiction.
 Now observe that if $H\in N_{\neg b}$, $[f]_{H}\neq [g]_{H}$ as 
 $ f\restriction q=g\restriction q$ holds for no $q\leq \neg b$. We conclude that $\dot{f}[N_{\neg b}]$ and $\dot{g}[N_{\neg b}]$ separate $\dot{f}(G)$ from $\dot{g}(G)$. 
% For the second point, let $\F\in\Presh(P)$ be strongly separated and assume $J$ to be subdense. We want to prove that $\Lambda^J_\F=\Lambda^{J^1_P}_\F$. To this end, it is enough to show that, if $[f]_{J,G}\neq [h]_{J,G}$, then $\{r\in P: f\restriction r=h\restriction r\}$ is not dense below any element $p\in G$. 
% But this is immediate, otherwise we would have $p\in G$ such that $f\restriction p=h\restriction p$ by the strong separativity of $\F$, which would give that $[f]_{J,G}= [h]_{J,G}$ as witnessed by $\downarrow p\in J(p)$.
% Now, Let $q,r\in G$ be such that $f\in \F(q), h\in \F(r)$. Then $p=q\wedge r\in G$ as well, and
%w.l.o.g. by renaming $f:=f\restriction p$, $h:=h\restriction p$ we can suppose $f,h\in\F(p)$.
%
%Consider the Boolean value 
%\[
%N_s=\bigvee\bp{N_r: r\leq p,\, f\restriction r=h\restriction r}.
%\]
%If $s\in G$, then $f\restriction s= h\restriction s$ (since $\F$ is strongly separated), yielding 
%$[f]_{J,G}=[h]_{J,G}$
% \end{proof}
% 
% 
% \begin{lemma}
%\label{lemma:charHaus}
%Let $\bool{B}$ be a  Boolean algebra  (possibly non-complete), and $\pi:E\to\St(\bool{B})$ be an \'etal\'e bundle on $\St(\bool{B})$ with $(E,\sigma)$ a topological space.
%
%The following are equivalent: %if $P$ is $\bool{B}^+$ for some Boolean algebra $\bool{B}$:
%\begin{itemize}
%\item
% $(E,\sigma)$ is Hausdorff;
% \item
% $(E,\sigma)$ is $0$-dimensional.
%\end{itemize}
%Finally if $(E,\sigma)$ is Hausdorff it is also locally compact.
%% if $\F$ is weakly $J$-separated as witnessed by $D$ dense subset of $P$,
%%$\dot{f}[N_p]$ is a proper section for all $p\in D$.
%\end{lemma}
%
%\begin{proof}


\item[$\Lambda^1_\F$ is locally compact] 
We note that for $s:N_b\to\Lambda_\F$ local section with open range, $s[N_b]$ is a compact and open subspace of $\Lambda_\F$, being it homeomorphic to the compact Hausdorff space $N_b$.
Being $\Lambda_\F$ Hausdorff, its compact subsets are closed, hence the topology of $\Lambda_\F$ has a base of clopen compact sets and is therefore $0$-dimensional and locally compact. 

\item[$\Lambda^1_\F$ is extremally disconnected] It suffices to prove that the closure of an open set is open.

It is convenient to adopt this piece of notation for $U\subseteq X\subseteq\Lambda_\F$: $\mathrm{Cl}_X(U)$ denotes the closure of $U$ as a subset of $X$ relative to the subspace topology inherited by $X$ from $\Lambda_\F$. Note that $\mathrm{Cl}_X(U)=\mathrm{Cl}_{\Lambda_\F}(U)\cap X$ and that  $\mathrm{Cl}_X(U)=\mathrm{Cl}_{\Lambda_\F}(U)$ when $X$ is a closed subset of $\Lambda_\F$.

Assume $A$ is open for $\Lambda_\F$ and $[f]_G\in\mathrm{Cl}_{\Lambda_\F}(A)$, we show that $[f]_G$ admits an open neighborhood totally contained in $\mathrm{Cl}_{\Lambda_\F}(A)$. 

This will give that $\mathrm{Cl}_{\Lambda_\F}(A)$ is open as well. 

Towards this aim note that for some $b\in G$ we have that $[f]_G\in \dot{f}[N_b]$ and that $\dot{f}$ implements an homeomorphism of $N_b$ with $\dot{f}[N_b]$ which is a clopen subspace of $\Lambda_\F$.

We therefore get that 
\[
Y=\mathrm{Cl}_{\dot{f}[N_b]}(A\cap \dot{f}[N_b])=\mathrm{Cl}_{\Lambda_\F}(A\cap \dot{f}[N_b])
\]
and that $[f]_G\in Y$.

Since $\dot{f}[N_b]$ with subspace topology is homeomorphic to $N_b$ -- which is extremally disconnected -- and $A\cap \dot{f}[N_b]$ is open in this subspace, we get that $Y$ is an open subset of $\dot{f}[N_b]$ in the subspace topology. Since $\dot{f}[N_b]$ is open in $\Lambda_\F$, we get that $Y$ is an open subset of $\Lambda_\F$. Clearly $Y$ is totally contained in $\mathrm{Cl}_{\Lambda_\F}(A)$. 
The proof is completed.
\end{description}
\end{proof}

\end{document}
